\section{Visual Synthesis: From Equation to Dynamics}
\label{sec:ch4-visual-synthesis}

The same dynamical law can be represented in several complementary ways. No single representation is sufficient in every problem.

\begin{figure}[htbp]
\centering
\begin{tikzpicture}[node distance=1.5cm and 1.7cm,every node/.style={font=\small,align=center}]
\node[draw,rounded corners,fill=softblue,minimum width=3.4cm,minimum height=1cm] (eq) {Equation\\$\dot{\mathbf{x}}=\mathbf{f}$};
\node[draw,rounded corners,fill=softgreen,right=of eq,minimum width=3.4cm,minimum height=1cm] (field) {Direction field\\local tendencies};
\node[draw,rounded corners,fill=softgold,below=of field,minimum width=3.4cm,minimum height=1cm] (phase) {Phase portrait\\all trajectories};
\node[draw,rounded corners,fill=softgray,left=of phase,minimum width=3.4cm,minimum height=1cm] (time) {Time series\\one experiment};
\draw[<->,thick] (eq)--(field);\draw[<->,thick] (field)--(phase);\draw[<->,thick] (phase)--(time);\draw[<->,thick] (time)--(eq);
\draw[<->,thick,dashed] (eq)--(phase);\draw[<->,thick,dashed] (field)--(time);
\end{tikzpicture}

\caption{Four complementary representations of an ODE: symbolic equation, direction field, time series, and phase portrait.}
\label{fig:ch4-representation-map}
\end{figure}

\begin{description}
\item[Equation.] Best for derivation, parameter dependence, and exact manipulation.
\item[Direction field.] Displays local tendencies without solving globally.
\item[Time series.] Shows what an observer records as time passes.
\item[Phase portrait.] Reveals the global organization of all initial conditions.
\end{description}

\begin{summarybox}[Core results of Commits 36--38]
\begin{itemize}
\item An ODE combines a local law of change with initial data to determine a trajectory.
\item Separable and linear first-order equations model decay, relaxation, and saturation.
\item Second-order linear equations classify free, damped, and driven oscillations through characteristic roots.
\item Higher-order equations become first-order systems, making matrix and eigenvalue methods available.
\item Phase space turns dynamics into geometry; equilibria and their stability are read from vector fields and linearization.
\item Nonlinear parameters can reorganize the set of equilibria through bifurcations.
\end{itemize}
\end{summarybox}

\begin{lookingaheadbox}[Looking ahead]
The next production stage will expand this core with a larger worked-example sequence, numerical integration in Python and Wolfram Language, driven and nonlinear applications, and a systematic bridge into mechanics and quantum evolution.
\end{lookingaheadbox}
